\documentclass[serif, mathsansserif, aspectratio=169]{beamer}   % [handout] to remove pause

\setbeamertemplate{title page}[default][left]  % left-align title page
\beamertemplatenavigationsymbolsempty % remove navigation symbols

\usepackage{sourcesanspro}
\usepackage[T1]{fontenc}
\usepackage[utf8]{inputenc}
\usepackage[document]{ragged2e}
\usepackage{hyperref}
\usepackage{fontawesome5}
\usepackage[tracking=true]{microtype}
\usepackage[scaled=0.90]{helvet}
\usepackage{caption}
\usepackage[scaled=0.85]{beramono}
\usepackage{tgpagella}
\usepackage{mathpazo}
\usepackage{mathtools}
\usepackage{amsmath}
\usepackage{amssymb}
\usepackage{amsthm}
\usepackage{scalerel}
\usepackage{bm}
\usepackage{bbm}
\usepackage{xcolor}
\usepackage{enumerate}
\usepackage[percent]{overpic}
\usepackage{soul}
\usepackage{tikz}
\usepackage{listings}
\usepackage{lstbayes}
\usepackage{booktabs}
\usepackage{bold-extra}
\usepackage{textpos}
\usepackage{multirow}
\usepackage{longtable}
\usepackage{tabularx}
\usepackage[flushleft]{threeparttable}
\usepackage{varwidth}

\usepackage{xpatch}
\usepackage[plain]{algorithm}
\usepackage{algorithmicx}
\usepackage[noend]{algpseudocode}

\usepackage{graphicx}
\usepackage{fancybox}


\usepackage[skins,theorems]{tcolorbox}
\tcbset{highlight math style={enhanced,colframe=red,colback=white,arc=0pt,boxrule=1pt}}

%%%%%%%%%%%%%%%%%%%%%%%%%%%%%%%%%%%%%%%%%%%%%%%%%%
% Bibliography settings
%%%%%%%%%%%%%%%%%%%%%%%%%%%%%%%%%%%%%%%%%%%%%%%%%%
\usepackage[backend=biber, style=apa, autocite=inline, doi=false, url=false]{biblatex}
\DeclareLanguageMapping{english}{english-apa}
\setcounter{biburlnumpenalty}{85}  % for urls that are too long
\addbibresource{bibliography.bib}
\setbeamertemplate{bibliography item}{}% Remove reference icon.
\setbeamercolor{bibliography entry author}{fg=bonnblue}
\setbeamercolor{bibliography entry note}{fg=bonnblue}
\renewcommand*{\bibfont}{\footnotesize}
\setlength\bibitemsep{\itemsep}
\AtEveryBibitem{%
  \clearfield{volume}%
  \clearfield{number}%
  \clearfield{pages}%
}
\newcommand\mycite[1]{\textcolor{bonnblue}{\citeauthor{#1}} (\textcolor{bonnblue}{\citeyear{#1}})}


%%%%%%%%%%%%%%%%%%%%%%%%%%%%%%%%%%%%%%%%%%%%%%%%%%
% Algorithm
%%%%%%%%%%%%%%%%%%%%%%%%%%%%%%%%%%%%%%%%%%%%%%%%%%
\makeatletter
\xpatchcmd{\algorithmic}{\itemsep\z@}{\itemsep=1ex plus1pt}{}{}
\makeatother



%%%%%%%%%%%%%%%%%%%%%%%%%%%%%%%%%%%%%%%%%%%%%%%%%%
% Math definitions
%%%%%%%%%%%%%%%%%%%%%%%%%%%%%%%%%%%%%%%%%%%%%%%%%%
\newcommand\indicator[1]{\scaleobj{1.4}{\mathbbm{1}} \! \left\{#1\right\}}
\newcommand\expectation[1]{\mathbb{E}\left[#1\right]}

\DeclareMathOperator*{\argmax}{argmax}


%%%%%%%%%%%%%%%%%%%%%%%%%%%%%%%%%%%%%%%%%%%%%%%%%%
% Colors
%%%%%%%%%%%%%%%%%%%%%%%%%%%%%%%%%%%%%%%%%%%%%%%%%%
\definecolor{bonnblue}{RGB}{4, 83, 156}  % hex: #04539C
\definecolor{bonngrey}{RGB}{148, 147, 132}  % hex: #8A9384
\definecolor{bonnyellow}{RGB}{251, 187, 6}  % hex: #FBBB06
\definecolor{bonngreen}{RGB}{0, 128, 0}  % hex: #FBBB06
\definecolor{bonnred}{RGB}{153, 0, 0}  % hex: #FBBB06
\definecolor{red}{RGB}{65, 16, 16}  % hex: #411010

%%%%%%%%%%%%%%%%%%%%%%%%%%%%%%%%%%%%%%%%%%%%%%%%%%
% Auxiliary files
%%%%%%%%%%%%%%%%%%%%%%%%%%%%%%%%%%%%%%%%%%%%%%%%%%
\newcommand{\email}[1]{\def\@email{\texttt{\MakeLowercase{\textls[10]{#1}}}}}

\setbeamerfont{title}{series=\bfseries} % family=\sourcesanspro
\setbeamerfont{subtitle}{series=\scshape, size*={8}{15}} % family=\sourcesanspro
\setbeamerfont{author}{family=\scshape}

\setbeamercolor{title}{fg=black}
\setbeamercolor{subtitle}{fg=bonngrey}
\setbeamercolor{author}{fg=bonnblue}
\setbeamercolor{date}{fg=bonnyellow}
\setbeamercolor{institute}{fg=bonngrey}
\setbeamercolor{email}{fg=bonnyellow}

\defbeamertemplate*{title page}{customized}[1][]
{
    \usebeamerfont{title}\usebeamercolor[fg]{title}\textls[100]{\MakeUppercase{\textbf{\inserttitle}}}\par
    \usebeamerfont{subtitle}\usebeamercolor[fg]{subtitle}\textls[100]{\textsc{\insertsubtitle}}\par\par
    {\usebeamercolor[fg]{date}\usebeamerfont{date}\insertdate\par}
    \vfill
    \usebeamerfont{author}\usebeamercolor[fg]{author}\textls[100]{\textsc{\insertauthor}}\par
    \usebeamerfont{institute}{\usebeamercolor[fg]{institute}\textls[100]{\textsc{\insertinstitute}}}\par
    \bigskip
}



%% Geometry

\geometry{left=1cm, right=1.5cm}


%%%%%%%%%%%%%%%%%%%%%%%%%%%%%%%%%%%%%%%%%%%%%%%%%%
% Auxiliary commands
%%%%%%%%%%%%%%%%%%%%%%%%%%%%%%%%%%%%%%%%%%%%%%%%%%
% \graphicspath{files/} % graphics path
\newcommand{\labelitem}{\tikz\draw[black,fill=bonnblue] (0,0) circle (.35ex); \,\,}

\newcommand{\blue}[1]{\textcolor{bonnblue}{#1}}
\newcommand{\yellow}[1]{\textcolor{bonnyellow}{#1}}
\newcommand{\grey}[1]{\textcolor{bonngrey}{#1}}
\newcommand{\red}[1]{\textcolor{red}{#1}}
\newcommand{\green}[1]{\textcolor{bonngreen}{#1}}
\def\ruleline#1{\hbox to \hsize{\hrulefill\quad\lower.3em\hbox{#1}\quad\hrulefill}}


\newcommand{\mycheck}{\textcolor{bonngreen}{\faIcon{check}}}
\newcommand{\mycross}{\textcolor{bonnred}{\faIcon{times}}}


%%%%%%%%%%%%%%%%%%%%%%%%%%%%%%%%%%%%%%%%%%%%%%%%%%
% Beamer settings
%%%%%%%%%%%%%%%%%%%%%%%%%%%%%%%%%%%%%%%%%%%%%%%%%%
\setbeamercolor{itemize item}{fg=black}
\setbeamercolor{itemize subitem}{fg=black}
\setbeamercolor{itemize subsubitem}{fg=black}
\setbeamercolor{section head}{fg=cardinal} % currently unused.
\setbeamercolor{footline}{fg=bonnblue}
\setbeamercolor{section in toc}{fg=black}

\setbeamertemplate{itemize item}[circle]
\setbeamertemplate{itemize subitem}{{\textendash}}
\setbeamertemplate{itemize subsubitem}[triangle]
\setbeamertemplate{section in toc}{{\inserttocsectionnumber.}~\inserttocsection}

\setbeamerfont{section in toc}{series=\scshape, size=\Large}
\setbeamerfont{frametitle}{series=\scshape} % \itshape
\setbeamercolor{frametitle}{fg=bonnblue}

\setbeamertemplate{frametitle}
{
    \vspace*{0.4cm}
    \insertframetitle \\ \footnotesize \textcolor{bonngrey}{\insertframesubtitle}
}

\AtBeginSection[]{
  \setbeamertemplate{footline}[frame number]{}
  \setbeamertemplate{navigation symbols}{}
  \setbeamertemplate{footline}{}
  \begin{frame}
  \vfill
  \centering
  \begin{beamercolorbox}[sep=8pt,center]{title}
      {\usebeamerfont{title}\usebeamercolor[fg]{title}{\ruleline{\textsc{\insertsectionhead}}}}\par%
  \end{beamercolorbox}
  \vfill
  \end{frame}
  \addtobeamertemplate{navigation symbols}{}{%
      \usebeamerfont{footline}%
      \usebeamercolor[fg]{author}%
      \hspace{1em}%
      \raisebox{2em}[0pt][0pt]{%
      \normalsize \insertframenumber \kern 0.25em
      }
  }
}

\AtBeginSubsection[]{
  \setbeamertemplate{footline}[frame number]{}
  \setbeamertemplate{navigation symbols}{}
  \setbeamertemplate{footline}{}
  \begin{frame}
  \vfill
  \centering
  \begin{beamercolorbox}[sep=8pt,center]{title}
      {\usebeamerfont{title}\usebeamercolor[fg]{title}{\textsc{{\normalsize
      \grey{\insertsectionhead}}\\[1em] \insertsubsectionhead}}}\par%
  \end{beamercolorbox}
  \vfill
  \end{frame}
  \addtobeamertemplate{navigation symbols}{}{%
      \usebeamerfont{footline}%
      \usebeamercolor[fg]{author}%
      \hspace{1em}%
      \raisebox{2em}[0pt][0pt]{%
      \normalsize \insertframenumber \kern 0.25em
      }
  }
}



\title{{\LARGE The Social Costs of AI}}
\subtitle{{\small Ethics, AI \& Economics}}
\author[Aimee van Wynsberghe \& Tim Mensinger \& Stefan H{\"o}se \& Markus Franke]
{%
   \texorpdfstring{
        \begin{columns}
            \column{.5\linewidth}
            \centering
            {{\large Aimee van Wynsberghe}\\\textcolor{bonngrey}{IWE}}
            \column{.5\linewidth}
            \centering
            {{\large Tim Mensinger}\\\textcolor{bonngrey}{Bonn Graduate School of Economics}}
        \end{columns}
        \vspace{0.5cm}
        \begin{columns}
            \column{.5\linewidth}
            \centering
            {{\large Stefan H{\"ose}}\\\textcolor{bonngrey}{DRZE}}
            \column{.5\linewidth}
            \centering
            {{\large Markus Franke}\\\textcolor{bonngrey}{IWE \& DRZE}}
        \end{columns}
   }
   {John Doe \& Jane Doe}
}
\email{}
\date{23. March 2022}


\begin{document}


\begin{frame}
    \maketitle
\end{frame}

% Page numbers
\addtobeamertemplate{navigation symbols}{}{%
    \usebeamerfont{footline}%
    \usebeamercolor[fg]{author}%
    \hspace{1em}%
    \raisebox{2em}[0pt][0pt]{%
    \normalsize \insertframenumber \kern 0.25em
    }
}
\begin{frame}
    \tableofcontents[hideallsubsections]
\end{frame}

\section{Who am I?}
\begin{frame}{Who am I?}
    \begin{itemize}
        \setlength\itemsep{1em}
        \item[\labelitem] From 2020: PhD Student in Economics, BGSE
        \item[\labelitem] Until 2020: Economics, Computer Science \& Math
        \item[\labelitem] Projects: Machine Learning in Economics
            \begin{itemize}
                \setlength\itemsep{0.5em}
                \item Modeling human decision making
                \item Predicting the effects of programs
            \end{itemize}
        \item[\labelitem] Research: Computational Econometrics
    \end{itemize}
\end{frame}


\section{The Research Question}
\begin{frame}{AI: An Example}
    \vspace{-1cm}
    Let's fix an example for this presentation
    \vspace{1cm}
    \begin{itemize}
        \setlength\itemsep{1em}
        \item[\labelitem] Language modeling (NLP)
        \item[\labelitem] Autocorrection on phones and computers
        \item[\labelitem] Amazon Alexa etc.
        \item[\labelitem] Google Translate etc.
    \end{itemize}
\end{frame}

\begin{frame}{Example contd.}
\begin{columns}
\begin{column}{0.5\textwidth}
\begin{overpic}[width=\textwidth]{files/venn}
 \put (30,68) {Artificial Intelligence}
 \put (20,40) {Machine Learning}
 \put (30,22) {\parbox{.5in}{Deep\\Learning}}
\end{overpic}
\end{column}
\begin{column}{0.5\textwidth}
    \begin{itemize}
        \setlength\itemsep{1em}
        \item[\labelitem] AI: Teaching computers languages
        \item[\labelitem] ML: Using specific models to do the above task
        \item[\labelitem] DL: Using ML models but with \textbf{many} \emph{parameters}
    \end{itemize}
\end{column}
\end{columns}
\end{frame}


\begin{frame}{What are parameters?}
\begin{itemize}
    \setlength\itemsep{1em}
    \item[\labelitem] Task: Write about your day
    \item[\labelitem] Catch: You can only use max. \grey{10}, \blue{100}, or \yellow{1000} words
    \item[\labelitem] Here: Number of words $\equiv$ Number of parameters
\end{itemize}
\end{frame}

\begin{frame}{The Problem}
    \begin{table}[]
    \renewcommand{\arraystretch}{2}
    \hskip-2.0cm
        \begin{tabular}{ll}
            \labelitem \textbf{Many parameters} & $\implies$ large computational costs\\
            \labelitem Large computational costs & $\implies$ large energy consumption\\
            \labelitem Large energy consumption & $\implies$ \textbf{large $CO_2$ emissions}
        \end{tabular}
    \end{table}
\end{frame}


\begin{frame}{Are the emission really that large?}
    \pause
    \textbf{Yes!}\\[1em]Training a large NLP model $\approx$ 125 round-trip
    flights New York - Beijing\\[0.5em]---\mycite{Strubell_2020}
\end{frame}

\begin{frame}{}
    \hspace{-1.5cm}
    \begin{centering}
    \begin{overpic}[width=1.1\textwidth]{files/model-parameters-timeline}
     \put (15, 15) {BERT Large: 340 Million}
     \put (80, 45) {T-NLG: 17 Billion}
     \put (20, 35) {$\approx$ 50 times as many parameters}
    \end{overpic}
    \end{centering}
\end{frame}

\begin{frame}{But what causes the emission?}
    \vspace{-1cm}
    For state-of-the-art large models:
    \vspace{1cm}
    \begin{itemize}
        \setlength\itemsep{1em}
        \item[\labelitem] \textbf{(Pre-) Training}\\
            $\approx$ 1 month on 1,000 computers; \mycite{Narayanan_2021}\\
        \item[\labelitem] \textbf{Fine-Tuning}\\
            $\approx$ few days on several computers; \mycite{Devlin_2018}\\
        \item[\labelitem] \textbf{Deployment}\\
            $\approx$ millisecond
    \end{itemize}
\end{frame}

\begin{frame}{How often do these stages happen?}
    \vspace{-1cm}
    Let's consider a language model used by Google:
    \vspace{1cm}
    \begin{itemize}
        \setlength\itemsep{1em}
        \item[\labelitem] \textbf{(Pre-) Training \blue{and Model Development}}\\
            $\approx$ 100 times (\emph{further resources needed})
        \item[\labelitem] \textbf{Fine-Tuning \blue{and Model Development}}\\
            $\approx$ 1,000 times (\emph{further resources needed})
        \item[\labelitem] \textbf{Deployment}\\
            $\approx$ 5 billion Google searches every day
    \end{itemize}
\end{frame}


\begin{frame}{The Problem contd.}
\begin{centering}
    \begin{tikzpicture}[fill opacity=0.2, text opacity=1, line width=0.1pt]

        \node (A) at (0, 8) {AI Model};
        \node (B) at (10, 8) {\textcolor{black}{$CO_2$}};

        \draw[line width=0.3mm, ->] (A) -- (B);

        \pause
        \node (C) at (5, 6) {Energy};
        \node (D) at (1, 3) {Hardware \& Software};
        \node (E) at (5, 3) {\parbox{1in}{{Pre-Training,} {Fine-Tuning,} Deployment}};

        \draw[line width=0.3mm, bonnblue, ->] (A) .. controls (0, 6.5) .. (C);
        \draw[line width=0.3mm, bonnblue, ->] (D) .. controls (3.5, 4) .. (C);
        \draw[line width=0.3mm, bonnblue, ->] (E) -- (C);

        \pause
        \node (F) at (9, 3) {Energy Source};
        \draw[line width=0.3mm, bonnyellow, ->] (C) .. controls (9, 6.5) .. (B);
        \draw[line width=0.3mm, bonnyellow, ->] (F) .. controls (11, 5) .. (B);

    \end{tikzpicture}
\end{centering}
\end{frame}

\begin{frame}{Research Questions}
    \begin{itemize}
        \setlength\itemsep{1em}
        \item[\labelitem] How do we quantify emissions of AI models?
        \item[\labelitem] Can we nudge the AI community towards sustainability?
    \end{itemize}
\end{frame}

\section{What is BERT and why do we need it?}

\begin{frame}{BERT}
    \begin{itemize}
        \setlength\itemsep{1em}
        \item[\labelitem] Language model developed at Google by \mycite{Devlin_2018}
        \item[\labelitem] Is used e.g. for Google search queries
        \item[\labelitem] Main example in \mycite{Strubell_2020} and \mycite{Schwartz_2020}
        \item[\labelitem] We use it, because
            \begin{itemize}
            \setlength\itemsep{0.5em}
                \item Perfect model size for our resources
                \item The literature is already familiar with the model
                \item We can compare our results to the literature
            \end{itemize}
    \end{itemize}
\end{frame}


\section{Project Status}

\begin{frame}{Project Status}
    \begin{itemize}
        \setlength\itemsep{1.5em}
        \item[\labelitem] Build test computing infrastructure \quad \mycheck
            \begin{itemize}
            \setlength\itemsep{0.7em}
                \item Setup energy measurement architecture \quad \mycheck
                \item Setup computing environment \quad \mycheck
                \item Run BERT model \quad \textcolor{bonnyellow}{\faIcon{spinner}}
            \end{itemize}
        \item[\labelitem] Build real computing infrastructure \quad \mycross
            \begin{itemize}
            \setlength\itemsep{0.7em}
                \item Get new computing hardware \quad \mycross
            \end{itemize}
        \item[\labelitem] Analyze energy data \quad \mycross
    \end{itemize}
\end{frame}


\begin{frame}[allowframebreaks]{References}
    \nocite{*}
    \printbibliography
\end{frame}

\setbeamertemplate{footline}[frame number]{}  % contact page has no frame number
\setbeamertemplate{navigation symbols}{}
\setbeamertemplate{footline}{}
\begin{frame}{Contact}

    \vspace{-1.5cm}
    \begin{picture}(320,250)
        \put(260,100){\includegraphics[scale=0.1]{files/qr-code-website}}
        \put(0,120){
            \begin{minipage}[t]{0.8\linewidth}{
                \begin{itemize}
                    \item[] \Large \faIcon{globe} \,\, \texttt{tmensinger.com}
                    \item[] \Large \faIcon{envelope} \, \texttt{tmensinger@uni-bonn.de}
                    \item[] \Large \faIcon{github} \,\, \texttt{github.com/timmens/social-cost-ai}
                \end{itemize}
            }
            \end{minipage}
        }
    \end{picture}

\end{frame}


\end{document}
